\section{Pour aller plus loin}

\begin{UPSTIexercice}{Rectangle vide}
    \begin{minipage}{0.8\textwidth}
        \UPSTIquestion{Écrire un programme qui affiche un rectangle vide de largeur \texttt{w} et de hauteur \texttt{h} données par l'utilisateur. Par exemple, pour \texttt{w=6} et \texttt{h=4}, le programme affichera le dessin ci-contre.}
    \end{minipage}%
    \hfill
    \begin{minipage}{0.1\textwidth}
        \begin{verbatim}
******
*    *
*    *
******
        \end{verbatim}
    \end{minipage}
\end{UPSTIexercice}

\begin{UPSTIprofOnlyEnv}
    \begin{UPSTIcorrectionP}{Rectangle vide}
        \begin{lstlisting}[language=C]
#include <stdio.h>
#include <cs50.h>
int main(void)
{
    int w, h;
    w = get_int("Entrez la largeur du rectangle : ");
    h = get_int("Entrez la hauteur du rectangle : ");

    for (int ligne = 1; ligne <= h; ligne = ligne + 1)
    {
        for (int col = 1; col <= w; col = col + 1)
        {
            if (ligne == 1 || ligne == h || col == 1 || col == w)
            {
                printf("*");
            }
            else
            {
                printf(" ");
            }
        }
        printf("\n");
    }
    return 0;
}
        \end{lstlisting}
    \end{UPSTIcorrectionP}
\end{UPSTIprofOnlyEnv}


\begin{UPSTIexercice}{Triangle de Pascal}
    \begin{minipage}{0.8\textwidth}
        Le triangle de Pascal est une disposition triangulaire des coefficients binomiaux. Chaque nombre est la somme des deux nombres directement au-dessus de lui dans le triangle.
        \UPSTIquestion{Écrire un programme qui affiche les \texttt{n} premières lignes du triangle de Pascal, où \texttt{n} est un entier donné par l'utilisateur. Par exemple, pour \texttt{n=5}, le programme affichera le dessin ci-contre.}
    \end{minipage}
    \hfill
    \begin{minipage}{0.15\textwidth}
        \begin{verbatim}
    1
   1 1
  1 2 1
 1 3 3 1
1 4 6 4 1
        \end{verbatim}
    \end{minipage}
\end{UPSTIexercice}

\begin{UPSTIprofOnlyEnv}
    \begin{UPSTIcorrectionP}{Triangle de Pascal}
        \begin{lstlisting}[language=C]
#include <stdio.h>
#include <cs50.h>

int main(void)
{
    int n;
    n = get_int("Nombre de lignes : ");
    for (int ligne = 0; ligne < n; ligne = ligne + 1)
    {
        int valeur = 1;
        for (int col = 0; col <= ligne; col = col + 1)
        {
            printf("%d ", valeur);
            valeur = valeur * (ligne - col) / (col + 1);
        }
        printf("\n");
    }
    return 0;
}
        \end{lstlisting}
    \end{UPSTIcorrectionP}
\end{UPSTIprofOnlyEnv}


\begin{UPSTIexercice}{Prédire l'affichage de triangles d'étoiles}
    \UPSTIquestion{Prédire l'affichage généré par le code suivant.
        \begin{lstlisting}[language=C]
#include <stdio.h>
int main(void) {
    int n = 5;
    for (int i = 1; i <= n; i++) {
        for (int j = 1; j <= i; j++) {
            printf("*");
        }
        printf("\n");
    }
    return 0;
}
        \end{lstlisting}}
    \UPSTIquestion{Prédire l'affichage généré par le code suivant.
        \begin{lstlisting}[language=C]
#include <stdio.h>
int main(void) {
    int n = 5;
    for (int i = n; i >= 1; i--) {
        for (int j = 1; j <= i; j++) {
            printf("*");
        }
        printf("\n");
    }
    return 0;
}
        \end{lstlisting}}
    \UPSTIquestion{Prédire l'affichage généré par le code suivant.
        \begin{lstlisting}[language=C]
#include <stdio.h>
int main(void) {
    int n = 5;
    for (int i = 1; i <= n; i++) {
        for (int j = 1; j <= n - i; j++) {
            printf(" ");
        }
        for (int k = 1; k <= i; k++) {
            printf("*");
        }
        printf("\n");
    }
    return 0;
}
        \end{lstlisting}}
\end{UPSTIexercice}

\begin{UPSTIprofOnlyEnv}
    \begin{UPSTIcorrectionP}{Prédire l'affichage de triangles d'étoiles}
        Premier exercice :
        \begin{verbatim}
*
**
***
****
*****
        \end{verbatim}
        Deuxième exercice :
        \begin{verbatim}
*****
****
***
**
*
        \end{verbatim}
        Troisième exercice :
        \begin{verbatim}
    *
   **
  ***
 ****
*****
        \end{verbatim}
    \end{UPSTIcorrectionP}
\end{UPSTIprofOnlyEnv}


\begin{UPSTIexercice}{Disque et cercle (Difficile)}
    \UPSTIquestion{Écrire un programme qui affiche un disque plein et un cercle vide de rayon \texttt{r} donné par l'utilisateur. Par exemple, pour \texttt{r=5}, le programme affichera :}
    \begin{verbatim}
  *****
 *******
*********
*********
*********
*********
 *******
  *****
    \end{verbatim}
\end{UPSTIexercice}

\begin{UPSTIprofOnlyEnv}
    \begin{UPSTIcorrectionP}{Très difficile : disque et cercle}
        Voici un exemple de solution pour afficher un cercle vide. L'affichage d'un disque plein est laissé en exercice.
        \begin{lstlisting}[language=C]
#include <stdio.h>
#include <math.h>
#define PI 3.14159265358979323846
#include <cs50.h>
int main(void) {
    int r;
    r = get_int("Entrez le rayon du cercle : ");

    for (int y = -r; y <= r; y++) {
        for (int x = -r; x <= r; x++) {
            if (fabs(sqrt(x*x + y*y) - r) < 0.5) {
                printf("*");
            } else {
                printf(" ");
            }
        }
        printf("\n");
    }
    return 0;
}
        \end{lstlisting}
    \end{UPSTIcorrectionP}
\end{UPSTIprofOnlyEnv}
